\documentclass[11pt]{article}
\usepackage[margin=1in]{geometry}
\usepackage[hidelinks]{hyperref}
\hypersetup{
  pdftitle={JDIQ Submission Highlights},
  pdfauthor={},
  pdfsubject={Anonymous submission highlights},
  pdfcreator={LaTeX},
  pdfproducer={LaTeX},
}

\pagestyle{empty}

\begin{document}

\begin{center}
{\Large Submission Highlights}\\[0.75em]
{\normalsize Data Quality for Derived Preference Graphs: Construction Sensitivity and Repair Outcomes in Multi-Ranker Retrieval}
\end{center}

\vspace{0.75em}

\begin{itemize}
  \item Preference graphs built from heterogeneous retrieval scores are treated as derived data artifacts, not as trusted raw inputs to graph repair.
  \item A compact seven-dimension data-quality audit framework separates provenance, calibration, vote semantics, conflict structure, repair quality, downstream utility, and reproducibility.
  \item Raw score margins are strongly scale dominated: BM25's conditional mean edge-weight share is 0.988 under raw construction but 0.512 after per-query, per-ranker normalization.
  \item Repair changes graph structure and, under genuine $P > k$ evaluation, can change top-$k$ membership, but no repaired-versus-unrepaired nDCG family in this study survives the primary Holm corrections.
  \item Exact SCIP repair improves the structural objective without revealing a reliable retrieval gain, yielding a practitioner-facing conclusion: normalization, vote construction, pooling, and evaluation cutoffs must be reported as first-class data-quality choices.
\end{itemize}

\end{document}
